
\subsection{Comparative Models and Regional Knowledge Ablations}
\label{subsec:comparative-models}

Three models, denoted as M0, M1, and M2, are considered for three-dimensional surface-pressure prediction. M0 uses only target-domain training data without incorporating pretrained aerodynamic knowledge. M1 introduces the sectional pressure prior and frozen features obtained from the infinite-span swept-wing model into the same three-dimensional prediction backbone. M2 further incorporates the wing-root and wing-tip pressure corrections and learns the target-domain residual relative to the complete hierarchical pressure prior. The three models share the same three-dimensional prediction backbone and differ primarily in the pretrained knowledge inputs and activated knowledge branches. The comparison between M0 and M1 evaluates the contribution of sectional knowledge transfer, whereas the comparison between M1 and M2 assesses the additional effect of regional pressure corrections.

To examine the separate contributions of wing-root and wing-tip knowledge, two additional ablation models, denoted as AR and AT, are considered at $N=20$. AR retains the sectional knowledge and wing-root pressure correction while disabling the wing-tip knowledge branch. Conversely, AT retains the sectional knowledge and wing-tip pressure correction while disabling the wing-root branch. For each ablation model, the pressure prior is recalibrated according to the activated knowledge components, and the target-domain prediction network is trained independently rather than disabling the corresponding branch of an already trained M2 model. The overall and spanwise regional pressure-prediction errors of M1, AR, AT, and M2 are compared to evaluate the individual and combined effects of the regional knowledge components.
